\subsection{the LHC}

In addition to its flagship high luminosity \pxp\ program the LHC also provides special runs aimed specifically for the field of heavy ion physics.  A list of these runs are compiled in Table~\ref{tab:lhc_nuclear_runs}.
These include heavy ion runs (A+A) to study QGP physics, small systems runs (\pxPb\ and $A_{\mathrm{light}}+A_{\mathrm{light}}$) to study cold nuclear matter effects and search for the turn on of QGP physics, and \pxp\ reference runs as a comparision point for the other systems at the same energy.
The author of this thesis was on site assiting with operations and data taking for all Run 3 run sans the 2026 \PbxPb\ run.

\begin{table}[!ht]
    \centering
    \begin{tabular}{| c | c | c | c | c |}
        \hline
        \textbf{Year(s)} & \textbf{Run} & \textbf{Energy ($\SNN$)} & \textbf{Species} & \textbf{Luminosity} \\
        \hline
        \hline
        2010 & Run 1 & 2.76 \TeV & \PbxPb & $\sim 10$ $\mu\text{b}^{-1}$ \\
        2011 & Run 1 & 2.76 \TeV & \PbxPb & $\sim 0.15$ \NB \\
        2013 & Run 1 & 2.76 \TeV & $pp$ (ref.) & $\sim 5$ pb$^{-1}$ \\
        2013 & Run 1 & 5.02 \TeV & \pxPb & $\sim 30$ \NB \\
        2015 & Run 2 & 5.02 \TeV & \PbxPb & $\sim 0.5$ \NB \\
        2015 & Run 2 & 5.02 \TeV & $pp$ (ref.) & $\sim 28$ pb$^{-1}$ \\
        2016 & Run 2 & 5.02 \TeV & \pxPb & $\sim 0.5$ \NB \\
        2016 & Run 2 & 8.16 \TeV & \pxPb & $160$ \NB \\
        2017 & Run 2 & 5.44 \TeV & Xe+Xe & $\sim 0.3$ $\mu\text{b}^{-1}$ \\
        2017 & Run 2 & 5.02 \TeV & $pp$ (ref.) & $\sim 270$ pb$^{-1}$ \\
        2018 & Run 2 & 5.02 \TeV & \PbxPb & $\sim 1.8$ \NB \\
        2023 & Run 3 & 5.36 \TeV & \PbxPb & $\sim 1.8$ \NB \\
        2023 & Run 3 & 5.36 \TeV & $pp$ (ref.) & $\sim 400$ pb$^{-1}$ \\
        2024 & Run 3 & 5.36 \TeV & \PbxPb & $\sim 1.7$ \NB \\
        2024 & Run 3 & 5.36 \TeV & \OxO & $\sim 8$ \NB \\
        2024 & Run 3 & 5.36 \TeV & \NexNe & $\sim 1$ \NB \\
        2025 & Run 3 & 5.36 \TeV & \PbxPb & $\sim 2.8$ \NB \\
        2026 & Run 3 & 5.36 \TeV & \PbxPb & $\sim 2.8$ \NB \\
        \hline
    \end{tabular}
    \caption{Summary of nuclear physics data-taking runs at the Large Hadron Collider (LHC), detailing runs for nuclear physics.}
    \label{tab:lhc_nuclear_runs}
\end{table}

\subsubsection{2016 \texorpdfstring{$p$}{p}+Pb Data}
In 2016, the LHC collided the asymmetric proton--Lead (\pxPb) system.
Two different center of mass energies were used, \SNN = 5.02 \TeV\ and \SNN = 8.16 \TeV.
The luminosity sampled by the higher energy sample far exceeded that of the lower energy sample.  Part of the reason for that discrepency comes from the difference in the average number of collisions per bunch crossing (\avmu) in each sample is much higher in the \SNN = 8.16 \TeV\ sample.
Higher \avmu\ also means higher level of contamination from in-time and out-of-time pileup events.

Another particular feature of the 2016 \pxPb\ dataset it fact that the direction of circulation for the proton and lead beams switched in the middle of the run.  In the first hi-\avmu\ period, \textbf{Period B}, protons circulated in the accelerator in a clockwise direction\footnote{direction defined from viewing the LHC from above} while lead ions circulated in the counterclockwise direction.  In the second hi-\avmu\ period, \textbf{Period C}, this directionality was reveresed.
\subsubsection{2025 Light Ion Data}


\subsubsection{Luminosity}
